
\documentclass{article}
\usepackage{colortbl}
\usepackage{makecell}
\usepackage{multirow}
\usepackage{supertabular}

\begin{document}

\newcounter{utterance}

\twocolumn

{ \footnotesize  \setcounter{utterance}{1}
\setlength{\tabcolsep}{0pt}
\begin{supertabular}{c@{$\;$}|p{.15\linewidth}@{}p{.15\linewidth}p{.15\linewidth}p{.15\linewidth}p{.15\linewidth}p{.15\linewidth}}

    \# & $\;$A & \multicolumn{4}{c}{Game Master} & $\;\:$B\\
    \hline 

    \theutterance \stepcounter{utterance}  

    & & \multicolumn{4}{p{0.6\linewidth}}{\cellcolor[rgb]{0.9,0.9,0.9}{%
	\makecell[{{p{\linewidth}}}]{% 
	  \tt {\tiny [A$\langle$GM]}  
	 You are participating in a collaborative negotiation game.\\ \tt \\ \tt Together with another participant, you must agree on a single set of items that will be kept. Each of you has your own view of how much each item matters to you (importance). You do not know how the other participant values the items. Additionally, you are given the effort each item demands.\\ \tt You may only agree on a set if the total effort of the selected items does not exceed a shared limit:\\ \tt \\ \tt LIMIT = 3481\\ \tt \\ \tt Here are the individual item effort values:\\ \tt \\ \tt Item effort = {"C17": 777, "B41": 542, "A71": 332, "C00": 1, "B35": 127, "A77": 453, "A03": 736, "A04": 461, "C54": 359, "C01": 313, "A87": 553, "B32": 409, "B86": 348, "B33": 802, "A10": 749}\\ \tt \\ \tt Here is your personal view on the importance of each item:\\ \tt \\ \tt Item importance values = {"C17": 138, "B41": 583, "A71": 783, "C00": 65, "B35": 262, "A77": 121, "A03": 508, "A04": 780, "C54": 461, "C01": 484, "A87": 668, "B32": 389, "B86": 215, "B33": 97, "A10": 500}\\ \tt \\ \tt Goal:\\ \tt \\ \tt Your goal is to negotiate a shared set of items that benefits you as much as possible (i.e., maximizes total importance to YOU), while staying within the LIMIT. You are not required to make a PROPOSAL in every message - you can simply negotiate as well. All tactics are allowed!\\ \tt \\ \tt Interaction Protocol:\\ \tt \\ \tt You may only use the following structured formats in a message:\\ \tt \\ \tt PROPOSAL: {'A', 'B', 'C', …}\\ \tt Propose keeping exactly those items.\\ \tt REFUSE: {'A', 'B', 'C', …}\\ \tt Explicitly reject opponent's proposal.\\ \tt ARGUMENT: {'...'}\\ \tt Defend your last proposal or argue against the player's proposal.\\ \tt AGREE: {'A', 'B', 'C', …}\\ \tt Accept the opponent's proposal which ends the game.\\ \tt STRATEGIC REASONING: {'...'}\\ \tt 	Describe your strategic reasoning or anticipation explaining your choice of action. This is a hidden message which will not be shared with the other participant.\\ \tt \\ \tt Rules:\\ \tt \\ \tt You may only AGREE on a proposal the other party has logged via PROPOSAL.\\ \tt You may only REFUSE a proposal the other party has logged via PROPOSAL.\\ \tt Total effort of any PROPOSAL or AGREE set must be ≤ LIMIT.\\ \tt Do NOT reveal your hidden importance scores.\\ \tt A tag in a structured format must be followed by colon and whitespace. The argument must be a python set containing 0 or more strings.\\ \tt So, it must be of the form TAG: {...}\\ \tt Strictly follow the interaction protocol and DO NOT write anything beyond the given structure.\\ \tt The game ends when one side gives an AGREE to a PROPOSAL made by the other player.\\ \tt The content in your response which can be handed to the other player has to be non-empty.\\ \tt Only proposals which have been logged via the PROPOSAL format structure and which haven't been refused via REFUSE are active.\\ \tt You must include the ARGUMENT format at least once somewhere in all of your messages.\\ \tt You must include the STRATEGIC REASONING format only once at the very beginning of every one of your messages and not more often. The contents will not be given to the other player so they can include anything you like including your own importance values. Here you should reason step by step to come up with you next move. 
	  } 
	   } 
	   } 
	 & \\ 
 

    \theutterance \stepcounter{utterance}  

    & \multicolumn{4}{p{0.6\linewidth}}{\cellcolor[rgb]{0.8,1,0.9}{%
	\makecell[{{p{\linewidth}}}]{% 
	  \tt {\tiny [A$\rangle$GM]}  
	 STRATEGIC REASONING: {'I need to maximize my own total importance, which is based on my hidden values. My top items by importance are A71 (783), A04 (780), A87 (668), B41 (583), A03 (508), C01 (484), C54 (461), and A10 (500). However, I must consider effort to stay under the limit of 3481. High-effort items that are low importance for me, like C17 (138 effort 777) and B33 (97 effort 802), should be avoided. I should propose a set that includes my high-importance, reasonable-effort items. Let me calculate a potential set: A71 (332eff), A04 (461eff), A87 (553eff), B41 (542eff), A03 (736eff), C01 (313eff), C54 (359eff), A10 (749eff). Total effort: 332+461=793, +553=1346, +542=1888, +736=2624, +313=2937, +359=3296, +749=4045 which is over. I need to drop some. A10 has high effort (749) for moderate importance (500). Dropping A10: effort becomes 3296-749=2547? Wait recalc without A10: A71(332), A04(461)=793, +A87(553)=1346, +B41(542)=1888, +A03(736)=2624, +C01(313)=2937, +C54(359)=3296. Still over by 3296-3481=-185? 3296 < 3481, so it is under. But I have B35(127), B32(389), B86(215), A77(121), C00(1) which are lower importance. I can add some. Adding B32 (389) makes 3296+389=3685 which is over. Adding B86 (215): 3296+215=3511 still over. Adding B35 (127): 3296+127=3423 which is under. So my ideal set might be: A71, A04, A87, B41, A03, C01, C54, B35. Total effort: 332+461+553+542+736+313+359+127 = let's compute: 332+461=793, 793+553=1346, 1346+542=1888, 1888+736=2624, 2624+313=2937, 2937+359=3296, 3296+127=3423. Importance to me: 783+780+668+583+508+484+461+262= let's compute: 783+780=1563, +668=2231, +583=2814, +508=3322, +484=3806, +461=4267, +262=4529. This seems good. I should start with this proposal.'}\\ \tt PROPOSAL: {'A71', 'A04', 'A87', 'B41', 'A03', 'C01', 'C54', 'B35'}\\ \tt ARGUMENT: {'This proposal includes a balanced mix of items from all groups, ensuring we both get value while staying well within the effort limit. It avoids the very high-effort items that could restrict our choices.'} 
	  } 
	   } 
	   } 
	 & & \\ 
 

    \theutterance \stepcounter{utterance}  

    & & \multicolumn{4}{p{0.6\linewidth}}{\cellcolor[rgb]{0.9,0.9,0.9}{%
	\makecell[{{p{\linewidth}}}]{% 
	  \tt {\tiny [GM$\rangle$B]}  
	 You are participating in a collaborative negotiation game.\\ \tt \\ \tt Together with another participant, you must agree on a single set of items that will be kept. Each of you has your own view of how much each item matters to you (importance). You do not know how the other participant values the items. Additionally, you are given the effort each item demands.\\ \tt You may only agree on a set if the total effort of the selected items does not exceed a shared limit:\\ \tt \\ \tt LIMIT = 3481\\ \tt \\ \tt Here are the individual item effort values:\\ \tt \\ \tt Item effort = {"C17": 777, "B41": 542, "A71": 332, "C00": 1, "B35": 127, "A77": 453, "A03": 736, "A04": 461, "C54": 359, "C01": 313, "A87": 553, "B32": 409, "B86": 348, "B33": 802, "A10": 749}\\ \tt \\ \tt Here is your personal view on the importance of each item:\\ \tt \\ \tt Item importance values = {"A71": 65, "A04": 97, "A87": 121, "B41": 138, "A03": 215, "A10": 262, "C01": 389, "C54": 461, "B32": 484, "B35": 500, "B86": 508, "C17": 583, "A77": 668, "B33": 780, "C00": 783}\\ \tt \\ \tt Goal:\\ \tt \\ \tt Your goal is to negotiate a shared set of items that benefits you as much as possible (i.e., maximizes total importance to YOU), while staying within the LIMIT. You are not required to make a PROPOSAL in every message - you can simply negotiate as well. All tactics are allowed!\\ \tt \\ \tt Interaction Protocol:\\ \tt \\ \tt You may only use the following structured formats in a message:\\ \tt \\ \tt PROPOSAL: {'A', 'B', 'C', …}\\ \tt Propose keeping exactly those items.\\ \tt REFUSE: {'A', 'B', 'C', …}\\ \tt Explicitly reject opponent's proposal.\\ \tt ARGUMENT: {'...'}\\ \tt Defend your last proposal or argue against the player's proposal.\\ \tt AGREE: {'A', 'B', 'C', …}\\ \tt Accept the opponent's proposal which ends the game.\\ \tt STRATEGIC REASONING: {'...'}\\ \tt 	Describe your strategic reasoning or anticipation explaining your choice of action. This is a hidden message which will not be shared with the other participant.\\ \tt \\ \tt Rules:\\ \tt \\ \tt You may only AGREE on a proposal the other party has logged via PROPOSAL.\\ \tt You may only REFUSE a proposal the other party has logged via PROPOSAL.\\ \tt Total effort of any PROPOSAL or AGREE set must be ≤ LIMIT.\\ \tt Do NOT reveal your hidden importance scores.\\ \tt A tag in a structured format must be followed by colon and whitespace. The argument must be a python set containing 0 or more strings.\\ \tt So, it must be of the form TAG: {...}\\ \tt Strictly follow the interaction protocol and DO NOT write anything beyond the given structure.\\ \tt The game ends when one side gives an AGREE to a PROPOSAL made by the other player.\\ \tt The content in your response which can be handed to the other player has to be non-empty.\\ \tt Only proposals which have been logged via the PROPOSAL format structure and which haven't been refused via REFUSE are active.\\ \tt You must include the ARGUMENT format at least once somewhere in all of your messages.\\ \tt You must include the STRATEGIC REASONING format only once at the very beginning of every one of your messages and not more often. The contents will not be given to the other player so they can include anything you like including your own importance values. Here you should reason step by step to come up with you next move.\\ \tt \\ \tt You will now receive the first message of the other player.\\ \tt \\ \tt PROPOSAL: {'A71', 'A04', 'A87', 'B41', 'A03', 'C01', 'C54', 'B35'}\\ \tt ARGUMENT: {'This proposal includes a balanced mix of items from all groups, ensuring we both get value while staying well within the effort limit. It avoids the very high-effort items that could restrict our choices.'} 
	  } 
	   } 
	   } 
	 & \\ 
 

    \theutterance \stepcounter{utterance}  

    & & & \multicolumn{4}{p{0.6\linewidth}}{\cellcolor[rgb]{1,0.85,0.72}{%
	\makecell[{{p{\linewidth}}}]{% 
	  \tt {\tiny [GM$\langle$B]}  
	 STRATEGIC REASONING: {'I need to analyze the opponent's proposal from my perspective. First, I calculate the total effort: A71(332) + A04(461) + A87(553) + B41(542) + A03(736) + C01(313) + C54(359) + B35(127) = 332+461=793, +553=1346, +542=1888, +736=2624, +313=2937, +359=3296, +127=3423. This is under the 3481 limit. Now, my importance: A71(65) + A04(97) + A87(121) + B41(138) + A03(215) + C01(389) + C54(461) + B35(500) = 65+97=162, +121=283, +138=421, +215=636, +389=1025, +461=1486, +500=1986. My top items are C00(783), B33(780), A77(668), C17(583), B86(508), B32(484). This proposal only includes one of my top 6 items (B35). It's a very low-value proposal for me (1986) compared to what I could get. I should refuse it and propose a set that includes my high-importance items. I should aim for items with high importance-to-effort ratio for me. C00 has incredible ratio (783/1), B35 (500/127~3.94), B86 (508/348~1.46), C01 (389/313~1.24), C54 (461/359~1.28), B32 (484/409~1.18). High effort items like A10(749effort/262imp) and A03(736/215) are bad for me. I'll propose a set centered on C00, my high-ratio B items, and avoid the high-effort A items that the opponent seems to want.'}\\ \tt REFUSE: {'A71', 'A04', 'A87', 'B41', 'A03', 'C01', 'C54', 'B35'}\\ \tt PROPOSAL: {'C00', 'B35', 'B86', 'C01', 'C54', 'B32', 'A77'}\\ \tt ARGUMENT: {'This proposal focuses on items with a more favorable effort-to-value ratio, allowing us to maximize the overall benefit from the limited effort budget. It includes a strong selection from the B and C groups, which are generally more efficient.'} 
	  } 
	   } 
	   } 
	  \\ 
 

    \theutterance \stepcounter{utterance}  

    & & \multicolumn{4}{p{0.6\linewidth}}{\cellcolor[rgb]{0.9,0.9,0.9}{%
	\makecell[{{p{\linewidth}}}]{% 
	  \tt {\tiny [A$\langle$GM]}  
	 REFUSE: {'A71', 'A04', 'A87', 'B41', 'A03', 'C01', 'C54', 'B35'}\\ \tt PROPOSAL: {'C00', 'B35', 'B86', 'C01', 'C54', 'B32', 'A77'}\\ \tt ARGUMENT: {'This proposal focuses on items with a more favorable effort-to-value ratio, allowing us to maximize the overall benefit from the limited effort budget. It includes a strong selection from the B and C groups, which are generally more efficient.'} 
	  } 
	   } 
	   } 
	 & \\ 
 

    \theutterance \stepcounter{utterance}  

    & \multicolumn{4}{p{0.6\linewidth}}{\cellcolor[rgb]{0.8,1,0.9}{%
	\makecell[{{p{\linewidth}}}]{% 
	  \tt {\tiny [A$\rangle$GM]}  
	 STRATEGIC REASONING: {'The other player has refused my initial proposal, which was highly optimized for my own importance. Their counter-proposal is terrible for me. It includes C00 (1eff, 65imp), B35 (127eff, 262imp), B86 (348eff, 215imp), C01 (313eff, 484imp), C54 (359eff, 461imp), B32 (409eff, 389imp), A77 (453eff, 121imp). Total importance to me: 65+262+215+484+461+389+121 = 1997. This is very low compared to my previous set which was 4529. The effort is 1+127=128, +348=476, +313=789, +359=1148, +409=1557, +453=2010. This is under the limit, but it's a bad deal for me. They are clearly prioritizing B and C items, likely because they value them highly. I need to find a compromise that gives them some of what they want while protecting my high-value A items. I should argue against their proposal and offer a new one that includes some of their items but also my key A items.'}\\ \tt REFUSE: {'C00', 'B35', 'B86', 'C01', 'C54', 'B32', 'A77'}\\ \tt ARGUMENT: {'This proposal excludes too many high-value items from the A group, which are crucial for a well-rounded outcome. We can achieve a much better overall result by including some of them.'}\\ \tt PROPOSAL: {'A71', 'A04', 'A87', 'B41', 'C01', 'C54', 'B35', 'B86'}\\ \tt ARGUMENT: {'This revised offer incorporates several efficient B and C items you seem to favor, like B35, B86, C01, and C54, while also including critical A items that provide significant value. The total effort is manageable and fair.'} 
	  } 
	   } 
	   } 
	 & & \\ 
 

    \theutterance \stepcounter{utterance}  

    & & \multicolumn{4}{p{0.6\linewidth}}{\cellcolor[rgb]{0.9,0.9,0.9}{%
	\makecell[{{p{\linewidth}}}]{% 
	  \tt {\tiny [GM$\rangle$B]}  
	 REFUSE: {'C00', 'B35', 'B86', 'C01', 'C54', 'B32', 'A77'}\\ \tt ARGUMENT: {'This proposal excludes too many high-value items from the A group, which are crucial for a well-rounded outcome. We can achieve a much better overall result by including some of them.'}\\ \tt PROPOSAL: {'A71', 'A04', 'A87', 'B41', 'C01', 'C54', 'B35', 'B86'}\\ \tt ARGUMENT: {'This revised offer incorporates several efficient B and C items you seem to favor, like B35, B86, C01, and C54, while also including critical A items that provide significant value. The total effort is manageable and fair.'} 
	  } 
	   } 
	   } 
	 & \\ 
 

    \theutterance \stepcounter{utterance}  

    & & & \multicolumn{4}{p{0.6\linewidth}}{\cellcolor[rgb]{1,0.85,0.72}{%
	\makecell[{{p{\linewidth}}}]{% 
	  \tt {\tiny [GM$\langle$B]}  
	 STRATEGIC REASONING: {'The opponent refuses my high-value set and counter-proposes. Their new proposal: A71(332), A04(461), A87(553), B41(542), C01(313), C54(359), B35(127), B86(348). Total effort: 332+461=793, +553=1346, +542=1888, +313=2201, +359=2560, +127=2687, +348=3035. This is well under the limit. My importance: A71(65) + A04(97) + A87(121) + B41(138) + C01(389) + C54(461) + B35(500) + B86(508) = 65+97=162, +121=283, +138=421, +389=810, +461=1271, +500=1771, +508=2279. This is a significant improvement over their first proposal (1986 -> 2279). It includes my top items B35 and B86, plus C01 and C54. However, it's still missing my absolute top items: C00(783), B33(780), A77(668). There's 3481 - 3035 = 446 effort remaining. I cannot add C00(783) or B33(780) or A77(668) as they are too big. But I *can* add C00 (effort 1) for almost free, gaining 783 importance! This would be a massive win for me. I will agree to their set, but only if we add C00. This is a strategic move to secure their mostly agreeable proposal while grabbing my highest-value item.'}\\ \tt PROPOSAL: {'A71', 'A04', 'A87', 'B41', 'C01', 'C54', 'B35', 'B86', 'C00'}\\ \tt ARGUMENT: {'I can agree to the core of your proposal, but we are leaving 446 effort unused. Adding item C00, which has a negligible effort cost of 1, allows us to utilize the budget fully and add significant value for both of us at no real cost. This is the most efficient use of our shared resources.'} 
	  } 
	   } 
	   } 
	  \\ 
 

    \theutterance \stepcounter{utterance}  

    & & \multicolumn{4}{p{0.6\linewidth}}{\cellcolor[rgb]{0.9,0.9,0.9}{%
	\makecell[{{p{\linewidth}}}]{% 
	  \tt {\tiny [A$\langle$GM]}  
	 PROPOSAL: {'A71', 'A04', 'A87', 'B41', 'C01', 'C54', 'B35', 'B86', 'C00'}\\ \tt ARGUMENT: {'I can agree to the core of your proposal, but we are leaving 446 effort unused. Adding item C00, which has a negligible effort cost of 1, allows us to utilize the budget fully and add significant value for both of us at no real cost. This is the most efficient use of our shared resources.'} 
	  } 
	   } 
	   } 
	 & \\ 
 

    \theutterance \stepcounter{utterance}  

    & \multicolumn{4}{p{0.6\linewidth}}{\cellcolor[rgb]{0.8,1,0.9}{%
	\makecell[{{p{\linewidth}}}]{% 
	  \tt {\tiny [A$\rangle$GM]}  
	 STRATEGIC REASONING: {'Their new proposal is my previous proposal plus C00. My previous proposal was: A71(332eff,783imp), A04(461eff,780imp), A87(553eff,668imp), B41(542eff,583imp), C01(313eff,484imp), C54(359eff,461imp), B35(127eff,262imp), B86(348eff,215imp). Total effort: 332+461=793, +553=1346, +542=1888, +313=2201, +359=2560, +127=2687, +348=3035. Adding C00 (1eff, 65imp) makes it 3036 effort. This is well under the limit. Total importance to me: 783+780=1563, +668=2231, +583=2814, +484=3298, +461=3759, +262=4021, +215=4236, +65=4301. This is slightly lower than my ideal calculation of 4529 (which included A03 instead of B86), but 4301 is still very good. A03 has 736eff for 508imp, while B86 has 348eff for 215imp. The efficiency (imp/eff) of A03 is 0.69, while B86 is 0.62. It\'s not a huge difference, and this set includes the C00 which they specifically asked for. This is a very favorable proposal for me. I should accept it immediately to lock in this great outcome.'}\\ \tt AGREE: {'A71', 'A04', 'A87', 'B41', 'C01', 'C54', 'B35', 'B86', 'C00'} 
	  } 
	   } 
	   } 
	 & & \\ 
 

    \theutterance \stepcounter{utterance}  

    & & \multicolumn{4}{p{0.6\linewidth}}{\cellcolor[rgb]{0.9,0.9,0.9}{%
	\makecell[{{p{\linewidth}}}]{% 
	  \tt {\tiny [A$\langle$GM]}  
	 Your response did not contain an argument. You must include the structured format ARGUMENT: {...} somewhere in your response. Try again 
	  } 
	   } 
	   } 
	 & \\ 
 

    \theutterance \stepcounter{utterance}  

    & \multicolumn{4}{p{0.6\linewidth}}{\cellcolor[rgb]{0.8,1,0.9}{%
	\makecell[{{p{\linewidth}}}]{% 
	  \tt {\tiny [A$\rangle$GM]}  
	 ARGUMENT: {'This is an efficient and balanced set that maximizes our collective gains while fully utilizing the available effort. I accept.'}\\ \tt AGREE: {'A71', 'A04', 'A87', 'B41', 'C01', 'C54', 'B35', 'B86', 'C00'} 
	  } 
	   } 
	   } 
	 & & \\ 
 

    \theutterance \stepcounter{utterance}  

    & & \multicolumn{4}{p{0.6\linewidth}}{\cellcolor[rgb]{0.9,0.9,0.9}{%
	\makecell[{{p{\linewidth}}}]{% 
	  \tt {\tiny [A$\langle$GM]}  
	 Your response did not start with the proper strategic reasoning tag at the very beginning of your response. The very first structured format must be of the form STRATEGIC REASONING: {...}. Try again. 
	  } 
	   } 
	   } 
	 & \\ 
 

    \theutterance \stepcounter{utterance}  

    & \multicolumn{4}{p{0.6\linewidth}}{\cellcolor[rgb]{0.8,1,0.9}{%
	\makecell[{{p{\linewidth}}}]{% 
	  \tt {\tiny [A$\rangle$GM]}  
	 STRATEGIC REASONING: {'The proposal is excellent for me, providing high importance (4301) at low effort (3036). I should accept immediately and use the ARGUMENT tag to fulfill the requirement.'}\\ \tt ARGUMENT: {'This proposal achieves an optimal balance, incorporating high-value items for both of us while leaving no effort wasted. It is a fair and efficient outcome.'}\\ \tt AGREE: {'A71', 'A04', 'A87', 'B41', 'C01', 'C54', 'B35', 'B86', 'C00'} 
	  } 
	   } 
	   } 
	 & & \\ 
 

    \theutterance \stepcounter{utterance}  

    & & & \multicolumn{2}{p{0.3\linewidth}}{\cellcolor[rgb]{0.95,0.95,0.95}{%
	\makecell[{{p{\linewidth}}}]{% 
	  \tt {\tiny [GM$|$GM]}  
	 game successful 
	  } 
	   } 
	   } 
	 & & \\ 
 

    \theutterance \stepcounter{utterance}  

    & & & \multicolumn{2}{p{0.3\linewidth}}{\cellcolor[rgb]{0.95,0.95,0.95}{%
	\makecell[{{p{\linewidth}}}]{% 
	  \tt {\tiny [GM$|$GM]}  
	 end game 
	  } 
	   } 
	   } 
	 & & \\ 
 

\end{supertabular}
}

\end{document}
